% !TEX program = xelatex
% Lernplatt - by Can Siebert
% Kurs 71 (Dig 42) - Tag 12 - Loesungsheft
% RPA skalieren ohne Kontrollverlust: Governance, Audit Trail, Exceptions, Security by Design, Change Management
%
% Self-contained xelatex source. Kein biblatex, keine externen Bilder.

\documentclass[11pt,a4paper,oneside]{article}

\usepackage{polyglossia}
\setdefaultlanguage{german}

\usepackage[a4paper,margin=2.5cm]{geometry}

\usepackage{fontspec}
\defaultfontfeatures{Ligatures=TeX}

\IfFontExistsTF{Charter}{%
  \setmainfont{Charter}%
}{%
  \IfFontExistsTF{Bitstream Charter}{%
    \setmainfont{Bitstream Charter}%
  }{%
    \setmainfont{TeX Gyre Schola}%
  }%
}

\IfFontExistsTF{Source Sans 3}{%
  \setsansfont{Source Sans 3}%
}{%
  \IfFontExistsTF{Source Sans Pro}{%
    \setsansfont{Source Sans Pro}%
  }{%
    \setsansfont{TeX Gyre Heros}%
  }%
}

\newcommand{\caveatfontfallback}{\itshape}
\IfFontExistsTF{Caveat}{%
  \newfontfamily\caveatfont{Caveat}[Scale=1.15]%
}{%
  \newcommand\caveatfont{\caveatfontfallback}%
}

\setmonofont{Latin Modern Mono}

\usepackage{microtype}
\linespread{1.15}

\usepackage{xcolor}
\definecolor{lpInk}{RGB}{20,28,38}
\definecolor{lpAccent}{RGB}{157,74,26}
\definecolor{lpAccentLight}{RGB}{246,228,212}
\definecolor{lpAmber}{RGB}{222,138,30}
\definecolor{lpAmberLight}{RGB}{255,243,217}
\definecolor{lpAmberBorder}{RGB}{196,118,18}
\definecolor{lpGray}{RGB}{96,104,114}
\definecolor{lpGrayLight}{RGB}{238,240,243}
\definecolor{lpGreen}{RGB}{34,120,72}
\definecolor{lpGreenLight}{RGB}{222,238,228}
\definecolor{lpGreenBorder}{RGB}{24,96,58}

\definecolor{lpL1}{RGB}{34,120,72}
\definecolor{lpL2}{RGB}{22,90,140}
\definecolor{lpL3}{RGB}{170,42,52}

\usepackage[most]{tcolorbox}
\tcbuselibrary{breakable,skins}

\newtcolorbox{infobox}[1][Hinweis]{%
  enhanced, breakable,
  colback=lpAccentLight,
  colframe=lpAccent,
  boxrule=0.6pt,
  arc=2pt,
  left=10pt,right=10pt,top=8pt,bottom=8pt,
  fonttitle=\sffamily\bfseries\color{white},
  coltitle=white,
  title={#1},
  attach boxed title to top left={xshift=8pt,yshift=-8pt},
  boxed title style={size=small, colback=lpAccent, colframe=lpAccent, boxrule=0.4pt, arc=1pt},
  top=14pt
}

\newtcolorbox{loesungbox}[1][Musterl\"osung]{%
  enhanced, breakable,
  colback=lpGreenLight,
  colframe=lpGreenBorder,
  boxrule=0.6pt,
  arc=2pt,
  left=10pt,right=10pt,top=8pt,bottom=8pt,
  fonttitle=\sffamily\bfseries\color{white},
  coltitle=white,
  title={#1},
  attach boxed title to top left={xshift=8pt,yshift=-8pt},
  boxed title style={size=small, colback=lpGreenBorder, colframe=lpGreenBorder, boxrule=0.4pt, arc=1pt},
  top=14pt
}

\usepackage{enumitem}
\setlist{itemsep=2pt, topsep=4pt, parsep=0pt}
\setlist[itemize,1]{label={\textbullet},leftmargin=1.6em}
\setlist[itemize,2]{label={\textendash},leftmargin=1.4em}
\setlist[enumerate,1]{label=\arabic*., leftmargin=1.8em}

\usepackage{tabularx}
\usepackage{array}
\usepackage{booktabs}
\usepackage{longtable}

\usepackage{fancyhdr}
\usepackage{lastpage}
\pagestyle{fancy}
\fancyhf{}
\renewcommand{\headrulewidth}{0.3pt}
\renewcommand{\footrulewidth}{0pt}
\fancyhead[L]{\footnotesize\sffamily\color{lpGray}L\"osungsheft \textperiodcentered{} Kurs 71 \textperiodcentered{} Tag 12}
\fancyhead[R]{\footnotesize\sffamily\color{lpGray}RPA \textendash{} Governance \& Skalierung}
\fancyfoot[C]{\footnotesize\sffamily\color{lpGray}\thepage{} / \pageref{LastPage}}

\fancypagestyle{plain}{%
  \fancyhf{}%
  \renewcommand{\headrulewidth}{0pt}%
  \fancyfoot[C]{\footnotesize\sffamily\color{lpGray}\thepage{} / \pageref{LastPage}}%
}

\usepackage[explicit]{titlesec}

\titleformat{\section}[block]
  {\sffamily\Large\bfseries\color{lpAccent}}
  {\thesection.}{0.6em}{#1}
  [{\vspace{2pt}\color{lpAccent}\titlerule[0.6pt]}]

\titleformat{\subsection}[block]
  {\sffamily\large\bfseries\color{lpInk}}
  {\thesubsection}{0.6em}{#1}

\titlespacing*{\section}{0pt}{14pt}{8pt}
\titlespacing*{\subsection}{0pt}{10pt}{4pt}

\usepackage[hidelinks,unicode]{hyperref}

\newcommand{\akzent}[1]{{\caveatfont\color{lpAmber}#1}}
\newcommand{\fachbegriff}[1]{\textbf{#1}}

\newcommand{\niveaubadge}[2]{%
  \tcbox[on line, colback=#1, colframe=#1, coltext=white,
         boxrule=0pt, arc=2pt, left=5pt,right=5pt,top=1pt,bottom=1pt,
         nobeforeafter]{\sffamily\bfseries\footnotesize #2}%
}
\newcommand{\nivLi}{\niveaubadge{lpL1}{L1\,\textbullet\,Wissen}}
\newcommand{\nivLii}{\niveaubadge{lpL2}{L2\,\textbullet\,Anwendung}}
\newcommand{\nivLiii}{\niveaubadge{lpL3}{L3\,\textbullet\,Management}}

\newcommand{\zeit}[1]{%
  \tcbox[on line, colback=lpGrayLight, colframe=lpGrayLight, coltext=lpInk,
         boxrule=0pt, arc=2pt, left=5pt,right=5pt,top=1pt,bottom=1pt,
         nobeforeafter]{\sffamily\footnotesize\textbf{$\sim$\,#1}}%
}

\newcommand{\aufgabenkopf}[4]{%
  \par\vspace{6pt}
  \noindent{\sffamily\Large\bfseries\color{lpAccent}Aufgabe #1 \textendash{} #2}\par
  \vspace{2pt}
  \noindent{\color{lpAccent}\rule{\linewidth}{0.6pt}}\par
  \vspace{4pt}
  \noindent #3 \hfill \zeit{#4}\par
  \vspace{6pt}
}

\newcommand{\ytquery}[1]{%
  \par\vspace{4pt}
  \noindent{\footnotesize\sffamily\color{lpGray}\textbf{YouTube-Suchquery:} \texttt{#1}}\par
}

% =========================================================================
\begin{document}

\begin{titlepage}
  \pagestyle{empty}
  \vspace*{1.5cm}
  \noindent{\sffamily\footnotesize\color{lpGray}LERNPLATT \textperiodcentered{} BY CAN SIEBERT}\\[2pt]
  \noindent{\color{lpAccent}\rule{\linewidth}{1.2pt}}\\[4pt]
  \noindent{\sffamily\footnotesize\color{lpGray}L\"OSUNGSHEFT \textperiodcentered{} MODUL 5 \textperiodcentered{} TAG 12 VON 16 \textperiodcentered{} KURS 71 (Dig 42) \textperiodcentered{} 15.\,Juni 2026}

  \vspace{3cm}
  {\sffamily\fontsize{30}{34}\selectfont\bfseries\color{lpInk}L\"osungsheft\\Tag 12\par}

  \vspace{0.7cm}
  {\sffamily\large\color{lpGray}Musterl\"osungen zu den elf Aufgaben\\des Aufgabenhefts \textendash{} RPA-Governance,\\Audit Trail und Security by Design.\par}

  \vspace{1.5cm}
  {\caveatfont\LARGE\color{lpGreen}Musterl\"osungen \textperiodcentered{} nicht die einzige Antwort}\par

  \vfill

  \noindent\begin{minipage}{0.55\linewidth}
    {\sffamily\footnotesize\color{lpGray}Hinweis:}\\
    {\sffamily\small\color{lpInk}Musterl\"osungen sind Diskussionsbasis,\\nicht die einzige korrekte L\"osung.}\\[10pt]
    {\sffamily\footnotesize\color{lpGray}Begleitend:}\\
    {\sffamily\small\color{lpInk}Aufgabenheft \& Lehrskript Tag 12}
  \end{minipage}\hfill
  \begin{minipage}{0.4\linewidth}
    \raggedleft
    {\sffamily\footnotesize\color{lpGray}Dozent:}\\
    {\sffamily\small\color{lpInk}Can Siebert}\\[10pt]
    {\sffamily\footnotesize\color{lpGray}Niveau-Mix:}\\
    {\sffamily\small\color{lpInk}3\,L1 \textperiodcentered{} 5\,L2 \textperiodcentered{} 3\,L3}
  \end{minipage}

  \vspace{0.5cm}
  \noindent{\color{lpAccent}\rule{\linewidth}{0.6pt}}
\end{titlepage}

\section*{Hinweise zu den Musterl\"osungen}
\thispagestyle{plain}

\noindent Die folgenden Musterl\"osungen sind \emph{eine} m\"ogliche, didaktisch nachvollziehbare L\"osung \textendash{} sie sind nicht die einzig richtige. Insbesondere auf L2 und L3 sind mehrere Begr\"undungswege legitim, solange sie:

\begin{itemize}
  \item den im Lehrskript eingef\"uhrten Begriffsapparat (Lifecycle, RACI, Exception-Typen, Security by Design) sauber nutzen,
  \item Annahmen explizit machen, statt sie zu verschleiern,
  \item Zielkonflikte (Speed vs.\ Compliance, Standardisierung vs.\ Akzeptanz) benennen, statt sie wegzudefinieren,
  \item zu einer klar formulierten Entscheidung f\"uhren \textendash{} mit Fr\"uhwarnsignalen.
\end{itemize}

\begin{infobox}[Nutzung im Coaching]
Vergleichen Sie Ihre eigene Antwort \emph{vor} der Musterl\"osung. Wo weicht Ihre Argumentation ab? Ist die Abweichung eine andere Annahme \textendash{} oder ein Denkfehler? Genau dort entsteht der Lerneffekt.
\end{infobox}

\clearpage

% =========================================================================
% L1 AUFGABEN
% =========================================================================
\section*{Niveau L1 \textendash{} Wissen \& Reproduktion}
\addcontentsline{toc}{section}{L1 -- Wissen}

% --- AUFGABE 1 -----------------------------------------------------------
\aufgabenkopf{1}{Bot-Lifecycle \& Governance-Rollen}{\nivLi}{8\,min}

\ytquery{RPA bot governance lifecycle phases process owner roles}

\begin{loesungbox}
\textbf{Zuordnung Phase \textrightarrow{} prim\"are Rolle:}

\smallskip
\begin{tabularx}{\linewidth}{@{}p{3.0cm}p{4.5cm}X@{}}
\toprule
\textbf{Phase} & \textbf{Prim\"are Rolle} & \textbf{Was passiert konkret} \\
\midrule
Idea       & Process Owner / Change Owner & Bedarf identifizieren, Business-Case-Skizze. \\
Assessment & Bot Owner + Process Owner    & Eignung pr\"ufen (Volumen, Regelbasiertheit, Stabilit\"at, Risiko); Go/No-Go-Vorschlag. \\
Design     & Developer + Reviewer         & Soll-Flow, Exception-Logik, Datenzugriffe entwerfen. \\
Build      & Developer                    & Implementierung gegen Standards (Naming, Logging). \\
Test       & Reviewer/QA                  & Funktional + nicht-funktional (Last, Recovery); Pass/Fail-Protokoll. \\
Deploy     & Operations + Security Owner  & Freigabe, Account-Setup, Vault, Monitoring scharf schalten. \\
Run        & Operations                   & Health, Queue, Erfolgsquote, Exceptions; MTTR steuern. \\
Improve    & Bot Owner + Developer        & Lessons Learned, Optimierung, Versions-Update. \\
Retire     & Process Owner + Security     & Zugriffe entziehen, Archivierung Audit-Trail, Doku abschlie{\ss}en. \\
\bottomrule
\end{tabularx}

\smallskip
\textbf{Assessment:} Entscheidung \emph{Bot ja/nein} \textendash{} produziert ein Fit-Profil (Volumen, Stabilit\"at, Risiko) und einen Erstent\-wurf des Business Case.\\
\textbf{Test:} Funktionale und nicht-funktionale Pr\"ufung (Last, Recovery, Exception-Pfade); Pass/Fail-Protokoll als Voraussetzung f\"ur UAT/Deploy.\\
\textbf{Retire:} Zugriffe entziehen, Audit Trail archivieren, Prozess-Doku abschlie{\ss}en \textendash{} \emph{nicht} der ``Bot l\"oschen''-Klick, sondern ein eigener Schritt mit Owner.

\smallskip
\textbf{Risiken bei ausgelassener Improve-Phase:}
\begin{itemize}
  \item \emph{Drift}: UI-\"Anderungen werden nicht nachgef\"uhrt, Exception Rate steigt unbemerkt; der Bot verursacht im Stillen mehr Aufwand als er spart.
  \item \emph{Wissensverlust}: Erkenntnisse aus Incidents flie{\ss}en nicht in Standards zur\"uck; andere Bots wiederholen denselben Fehler.
\end{itemize}
\end{loesungbox}

\clearpage

% --- AUFGABE 2 -----------------------------------------------------------
\aufgabenkopf{2}{Vier nicht verhandelbare Governance-Standards}{\nivLi}{8\,min}

\ytquery{RPA governance standards naming versioning logging exception codes}

\begin{loesungbox}
\textbf{Vier Standards \textendash{} konkret und mit Konsequenz:}

\begin{enumerate}
  \item \textbf{Naming \& Versionierung.} Format: \texttt{<Domain>\_<Prozess>\_<Action>\_v<Major>.<Minor>}, z.\,B.\ \texttt{Finance\_Invoice\_Postin\-g\_v2.3}. Versionsspr\"unge nach SemVer: Major bei API-Bruch, Minor bei Erweiterung, Patch bei Bugfix. \emph{Konsequenz bei Verletzung:} Bot wird im Repository abgelehnt; ohne Versionsnummer kein Audit-Trail-Eintrag m\"oglich.
  \item \textbf{Logging-Minimum.} Pro Bot-Run gelogged: \texttt{BotID}, \texttt{Version}, \texttt{InputRef}, \texttt{Action}, \texttt{Target\-System}, \texttt{Result\-Status}, \texttt{ExceptionCode}, \texttt{ProcessedBy}, \texttt{Timestamp}. Sensible Felder maskiert oder gehasht. \emph{Konsequenz bei Verletzung:} Bot bleibt in Dev oder Test \textendash{} keine Prod-Freigabe.
  \item \textbf{Exception Codes.} Zentral gepflegtes Code-Verzeichnis pro Dom\"ane (\texttt{ERR\_INV\_001} = \emph{MissingCostCenter}, \texttt{ERR\_INV\_002} = \emph{DuplicateInvoice}\,\ldots). Codes referenzieren Playbooks. \emph{Konsequenz:} Ohne Exception-Code-Mapping wird ein neuer Bot nicht ausgerollt.
  \item \textbf{Freigabe vor Produktion.} Vier-Augen-Prinzip: Reviewer (Code/Standards) und Process Owner (fachliche Korrektheit), bei High-Risk-Bots zus\"atzlich Security Owner. Schriftliche Freigabe im Ticket. \emph{Konsequenz:} Bots ohne Freigabe-Ticket werden aus Prod entfernt \textendash{} hartes Stop-Kriterium.
\end{enumerate}

\smallskip
\textbf{F\"unfter Standard (Vorschlag):} \emph{Rollback-Plan vor Go-Live} \textendash{} jeder Bot ben\"otigt einen dokumentierten Stop- und Rollback-Pfad (alte Version, manuelle Notfall-Variante, Verantwortlicher) bevor er live geht. Begr\"undung: Sobald ein Bot l\"auft, geht jede Sekunde Stillstand auf Kosten des Fachbereichs \textendash{} Rollback muss \emph{vorher} klar sein, nicht im Incident verhandelt werden.
\end{loesungbox}

\clearpage

% --- AUFGABE 3 -----------------------------------------------------------
\aufgabenkopf{3}{Exception-Typen \& Audit-Trail-Felder}{\nivLi}{10\,min}

\ytquery{RPA audit trail exception handling business system data exception}

\begin{loesungbox}
\textbf{1. Exception-Typen mit drei Beispielen je Dom\"ane:}

\smallskip
\begin{tabularx}{\linewidth}{@{}p{2.6cm}XXX@{}}
\toprule
\textbf{Typ} & \textbf{Rechnungsverarbeitung} & \textbf{HR-Onboarding} & \textbf{Bestellabwicklung} \\
\midrule
Business   & ApprovalRequired ($>$\,10\,k\,\euro), DuplicateInvoice, ApprovalDeadlineMissed & ContractNotSigned, ManagerApprovalMissing, RoleNotDefined & CreditLimitExceeded, BlockedCustomer, MissingApproval \\
\midrule
System     & ERPDown, Timeout, OCR-Service down & AD-NotReachable, MailServerDown, HR-System Timeout & SAP-Down, WMS-Timeout, EDI-Channel down \\
\midrule
Data       & MissingCostCenter, UnknownVendor, WrongCurrency & MissingTaxID, InvalidIBAN, MissingDepartment & MissingArticleNumber, InvalidPriceList, MissingShippingAddress \\
\bottomrule
\end{tabularx}

\smallskip
\textbf{2. Audit-Log-Template (Mindestumfang):}

\begin{itemize}
  \item \texttt{LogID} (UUID, append-only)
  \item \texttt{BotID} + \texttt{Version} (z.\,B.\ \texttt{Finance\_Invoice\_Posting\_v2.3})
  \item \texttt{Timestamp} (ISO 8601, UTC, append-only)
  \item \texttt{InputRef} (z.\,B.\ Rechnungsnummer + Ticket-ID)
  \item \texttt{Action} (z.\,B.\ \emph{POST\_INVOICE}, \emph{UPDATE\_VENDOR})
  \item \texttt{TargetSystem} + \texttt{TargetObject}
  \item \texttt{ResultStatus} (\emph{Success} / \emph{BusinessException} / \emph{SystemException} / \emph{DataException}) \textendash{} append-only
  \item \texttt{ExceptionCode} (Referenz zum Playbook)
  \item \texttt{ProcessedBy} (Bot-Account; bei Handover: Name des Operators)
  \item \texttt{Reason} (Freitext, optional)
\end{itemize}

\smallskip
\textbf{3. Append-only-Felder:} \texttt{LogID}, \texttt{Timestamp}, \texttt{ResultStatus}. Begr\"undung: Ohne diese drei kann nachtr\"agliche Manipulation Schaden vert\"uschen \textendash{} Auditor verlangt \emph{Integrit\"at}, nicht nur \emph{Vollst\"andigkeit}. \"Anderungen werden \"uber neue Log-Eintr\"age dargestellt (Event-Sourcing-Prinzip).
\end{loesungbox}

\clearpage

% =========================================================================
% L2 AUFGABEN
% =========================================================================
\section*{Niveau L2 \textendash{} Anwendung \& Transfer}
\addcontentsline{toc}{section}{L2 -- Anwendung}

% --- AUFGABE 4 -----------------------------------------------------------
\aufgabenkopf{4}{Governance-Minimum in 45 Minuten}{\nivLii}{30\,min}

\ytquery{RPA minimum viable governance RACI release flow fast track}

\begin{loesungbox}
\textbf{1. RACI (Auszug, sechs Rollen \"uber f\"unf T\"atigkeiten):}

\smallskip
\begin{tabularx}{\linewidth}{@{}p{3.9cm}cccccc@{}}
\toprule
\textbf{T\"atigkeit} & \textbf{PO} & \textbf{Bot-Own.} & \textbf{Dev} & \textbf{QA} & \textbf{Sec} & \textbf{Ops} \\
\midrule
Bot-Idee bewerten     & R/A & C   & I   & I   & C   & I   \\
Code reviewen         & I   & C   & R   & A   & C   & I   \\
Freigabe vor Prod     & C   & R   & I   & C   & A   & C   \\
Run \& Monitoring     & I   & A   & I   & I   & I   & R   \\
Incident bearbeiten   & C   & A   & C   & I   & C   & R   \\
\bottomrule
\end{tabularx}

\smallskip
\textbf{2. Vier nicht verhandelbare Standards:}
\begin{itemize}
  \item \emph{Versionierung} (SemVer, Major-Minor-Patch im Bot-Namen).
  \item \emph{Logging-Minimum} (Felder wie in Aufgabe 3; sensible Daten maskiert).
  \item \emph{Exception Codes} (zentrales Verzeichnis, Mapping zu Playbooks).
  \item \emph{Freigabe vor Prod} (Vier-Augen, schriftlich im Ticket).
\end{itemize}

\smallskip
\textbf{3. Release Flow mit zwei Gates:}\\
\emph{Build} \textrightarrow{} \emph{Test} \textrightarrow{} \textbf{Gate 1} \emph{(QA + Standards-Check)} \textrightarrow{} \emph{UAT} \textrightarrow{} \textbf{Gate 2} \emph{(Process Owner + Security Owner)} \textrightarrow{} \emph{Deploy} \textrightarrow{} \emph{Monitoring}.\\
\emph{Gate 1} entscheidet \"uber technische Reife (Code-Standards, Tests, Logging). \emph{Gate 2} entscheidet \"uber fachliche und Risiko-Freigabe \textendash{} ist nicht delegierbar.

\smallskip
\textbf{4. Fast-Track-Regel:} Kleine Bots d\"urfen Gate 2 ohne Security Owner passieren, wenn \emph{alle} folgenden Bedingungen erf\"ullt sind:
\begin{itemize}
  \item kein Schreibzugriff auf Finance- oder HR-Systeme,
  \item kein Umgang mit PII oder Gesundheitsdaten,
  \item maximal 100 Transaktionen pro Tag,
  \item Read-only-Operation auf nicht-kritischen Systemen,
  \item dokumentierter Rollback-Pfad im Repository.
\end{itemize}
\emph{Guardrail:} Verletzt ein Fast-Track-Bot eine der vier Standards (Logging, Versionierung, Exception Codes, Freigabe), wird er sofort \emph{ohne Diskussion} aus Prod gezogen \textendash{} keine Ausnahmen.

\smallskip
\emph{Andere RACI-Auspr\"agungen sind legitim, solange Accountability eindeutig bleibt und Standards plus Gates klar sind.}
\end{loesungbox}

\clearpage

% --- AUFGABE 5 -----------------------------------------------------------
\aufgabenkopf{5}{Audit Trail Standard \& f\"unf Exception-Playbooks}{\nivLii}{25\,min}

\ytquery{RPA audit log retention exception playbook MTTR exception rate}

\begin{loesungbox}
\textbf{1. Audit-Log-Template (kompakt):}

\smallskip
\begin{tabularx}{\linewidth}{@{}p{3.6cm}p{2.0cm}p{1.4cm}X@{}}
\toprule
\textbf{Feld} & \textbf{Typ} & \textbf{Pflicht?} & \textbf{Anmerkung} \\
\midrule
\texttt{LogID}         & UUID      & ja & append-only \\
\texttt{Timestamp}     & ISO 8601  & ja & append-only \\
\texttt{BotID}         & String    & ja & inkl.\ Version \\
\texttt{Version}       & SemVer    & ja & \\
\texttt{InputRef}      & String    & ja & z.\,B.\ Rechnungsnummer \\
\texttt{Action}        & Enum      & ja & POST/UPDATE/READ \\
\texttt{TargetSystem}  & String    & ja & ERP, CRM, Portal \\
\texttt{ResultStatus}  & Enum      & ja & append-only \\
\texttt{ExceptionCode} & String    & nein & nur bei Exception \\
\texttt{ProcessedBy}   & String    & ja & Bot-Account / Operator \\
\texttt{Reason}        & Text      & nein & Begr\"undung \\
\bottomrule
\end{tabularx}

\smallskip
\textbf{Aufbewahrung:} Finance-Logs \emph{5 Jahre} (steuerlich begr\"undet, GoBD); Logs zentral, role-based Access (Bot Owner + Audit + Security), Export-API f\"ur Auditor.

\smallskip
\textbf{2. Playbooks (kompakt):}

\smallskip
\begin{tabularx}{\linewidth}{@{}p{3.4cm}p{2.4cm}Xp{1.4cm}p{2.6cm}@{}}
\toprule
\textbf{Exception} & \textbf{Owner} & \textbf{Reaktion} & \textbf{SLA} & \textbf{Eskalation} \\
\midrule
MissingCostCenter   & Buchhaltung    & Queue f\"ur Klassifikation; Mail an Anforderer        & 24\,h & Teamleitung nach 48\,h \\
DuplicateInvoice    & Buchhaltung    & Stopp Bot; Ticket; Vorgang an Bearbeitenden          & 8\,h  & Process Owner nach 24\,h \\
ERPNotReachable     & IT-Operations  & Retry 3x exponentiell; Ticket an IT; Queue parken    & 2\,h  & Plattformbetrieb nach 2\,h \\
ApprovalThreshold   & Bot Owner      & Vier-Augen-Freigabe-Task an Genehmiger              & 24\,h & PO + n\"achste Hierarchie nach 48\,h \\
PortalUIChanged     & Bot Owner + Dev& Bot pausieren; Selektor-Update; Regression-Test     & 24\,h & Bereichsleitung nach 72\,h \\
\bottomrule
\end{tabularx}

\smallskip
\textbf{3. Zwei KPIs mit Anti-Gaming:}
\begin{itemize}
  \item \emph{Exception Rate} $\leq$\,12\,\%. \textbf{Anti-Gaming:} Z\"ahlung erfolgt \emph{automatisch} aus dem Audit-Log; manuelles ``Umlabeln'' einer Exception auf \emph{Success} l\"ost zus\"atzliche Eintragspflicht im Log aus (Reviewer + Begr\"undung).
  \item \emph{MTTR} $<$\,4\,h (von Exception bis L\"osung/Handover). \textbf{Anti-Gaming:} Stoppt nicht beim ``Ticket-erstellt'', sondern beim \emph{Resolved-with-Outcome}; Ticket-Schlie{\ss}ung ohne Resolution-Code wird abgewiesen.
\end{itemize}
\end{loesungbox}

\clearpage

% --- AUFGABE 6 -----------------------------------------------------------
\aufgabenkopf{6}{Security-by-Design Guardrails entwerfen}{\nivLii}{25\,min}

\ytquery{RPA security by design least privilege SoD credential vault bot account}

\begin{loesungbox}
\textbf{1. Account-Konzept:}
\begin{itemize}
  \item \emph{Bot-Accounts} sind technische Service-Accounts (\texttt{srv\_rpa\_invoice\_prod}, \texttt{srv\_rpa\_invoice\_dev}). Nie identisch mit Personal-Accounts; nie ``geteilt''.
  \item Passw\"orter und API-Keys liegen im \emph{Credential Vault} (z.\,B.\ konzeptionell HashiCorp Vault, CyberArk, Azure Key Vault). Bot fordert Secret zur Laufzeit an \textendash{} \emph{nie} im Code, nie in Konfigurationsdateien.
  \item Rotation: mindestens quartalsweise; bei Vorfall sofort.
\end{itemize}

\smallskip
\textbf{2. Least Privilege (drei Beispiel-Bots):}

\smallskip
\begin{tabularx}{\linewidth}{@{}p{3.0cm}X@{}}
\toprule
\textbf{Bot} & \textbf{Minimal-Berechtigungen} \\
\midrule
InvoiceBot     & ERP: \emph{Read} Kreditoren, \emph{Write} Buchungsbelege (nur ein Belegtyp); \emph{kein} Zugriff auf HR, kein Stammdaten-\"Andern. \\
HR-OnboardingBot& AD: nur \emph{User-Create} in zwei OUs; Mail: \emph{Send-as} auf eine Funktions-Mailbox; \emph{kein} Read auf andere Mitarbeiterdaten. \\
ReportingBot   & DWH: \emph{Read-only} auf zwei Views; Mail-Versand an feste Empf\"angerliste; \emph{kein} Schreibzugriff irgendwo. \\
\bottomrule
\end{tabularx}

\smallskip
\textbf{3. Drei SoD-Regeln (Finance-Bot):}
\begin{itemize}
  \item \emph{Anlegen} eines Kreditors und \emph{Freigeben} desselben Kreditors d\"urfen nicht durch denselben Account erfolgen.
  \item \emph{Erfassen} und \emph{Freigeben} einer Zahlung sind getrennt; Bot darf maximal eine der beiden Operationen.
  \item \emph{\"Anderung} von Stammdaten (Bankverbindung, Adresse) und \emph{Ausl\"osen} einer Zahlung d\"urfen nicht in derselben Bot-Session erfolgen (Cooling-off-Periode oder Vier-Augen).
\end{itemize}

\smallskip
\textbf{4. PII-/Datenschutz-Logik:}
\begin{itemize}
  \item \emph{Nie im Klartext geloggt:} Gehalt, Gesundheitsdaten, vollst\"andige IBAN, Geburtsdatum, Personalnummer.
  \item \emph{L\"osungen:} Hashing (z.\,B.\ Personalnummer als SHA-256), Maskierung (\texttt{DE89****6789}), oder ganz weglassen und Verweis auf ID nutzen.
  \item Zugriff auf PII-haltige Logs nur f\"ur Bot Owner + Audit + Datenschutzbeauftragter; Vier-Augen-Zugriff dokumentiert.
\end{itemize}

\smallskip
\textbf{5. Drei nicht verhandelbare Guardrails vor Go-Live:}
\begin{itemize}
  \item Bot-Account ist nicht identisch mit einem Personal-Account, Secrets im Vault.
  \item SoD-Pr\"ufung dokumentiert; kein Bot mit \emph{Anlegen + Freigeben} im selben Account.
  \item Logging-Integrit\"at (append-only) und PII-Maskierung gepr\"uft \textendash{} Security Sign-off im Ticket.
\end{itemize}
\end{loesungbox}

\clearpage

% --- AUFGABE 7 -----------------------------------------------------------
\aufgabenkopf{7}{Change-Paket f\"ur Mitarbeitende ``hinter dem Bot''}{\nivLii}{20\,min}

\ytquery{RPA change management role shift training resistance bot operator owner}

\begin{loesungbox}
\textbf{1. Rollenwandel:} Die T\"atigkeit verschiebt sich von \emph{repetitiver Eingabe} (Daten erfassen, Standardf\"alle abarbeiten) zu \emph{Steuerung} (Monitoring der Bot-Performance, Ausnahmen bewerten, fachliche Eskalation). Konkret: Was fr\"uher 80\,\% der Zeit ``Tippen'' war, wird zu 60\,\% ``Pr\"ufen / Entscheiden / Eskalieren''. Neue Skills: Pattern-Erkennung in Logs, kommunikative St\"arke (Lieferant ansprechen), grundlegendes Verst\"andnis von Bot-Status und Exception-Codes.

\smallskip
\textbf{2. Kommunikations-Botschaft (zwei S\"atze):}
``Der \emph{InvoiceBot} \"ubernimmt die Standardf\"alle, damit Sie mehr Zeit f\"ur die Ausnahmen, Lieferantenbeziehungen und Verbesserungen haben. Ihre Rolle wandelt sich von \emph{Erfasser} zu \emph{Entscheider} \textendash{} mit Training und Begleitung, und wir kommunizieren transparent, wenn sich an Stellenanzahl etwas \"andert.''

\smallskip
\textbf{3. Trainingspfad (drei Stufen):}

\smallskip
\begin{tabularx}{\linewidth}{@{}p{2.4cm}XX@{}}
\toprule
\textbf{Stufe} & \textbf{Inhalte (2)} & \textbf{Wirksamkeitsnachweis} \\
\midrule
Operator  & Bot-Status lesen; Standard-Exception aufl\"osen        & Aufl\"osung von 3 Standard-Exceptions selbstst\"andig in einer Woche \\
Owner     & Playbook-Pflege; Eskalations\-entscheidung treffen      & Eine eigenh\"andig dokumentierte Playbook-\"Anderung pro Quartal \\
Champion  & Verbesserungs\-vorschl\"age; Change-Multiplikator       & Eine umgesetzte Verbesserung pro Halbjahr; Mentor f\"ur 2 Operator \\
\bottomrule
\end{tabularx}

\smallskip
\textbf{4. Umgang mit aktivem Widerstand (``Ich werde diesen Bot nicht trainieren''):}
\begin{itemize}
  \item \emph{Zuh\"oren:} Vier-Augen-Gespr\"ach, offene Frage: ``Was l\"asst Sie das so sehen?'' \textendash{} Interessen / \"Angste / Verlustsorgen werden konkret. \emph{Zeit: 20\,min, nicht eilig.}
  \item \emph{Reframe:} ``Sie kennen den Prozess am besten \textendash{} ohne Sie wird der Bot Schrott. Wir brauchen Ihre Erfahrung, um die Ausnahmen zu modellieren.'' Verantwortung sichtbar machen, nicht moralisch.
  \item \emph{Klare Konsequenz:} ``Ohne Ihren Input geht der Bot trotzdem live \textendash{} aber schlechter, und Sie tragen die Folgen sp\"ater mit. Sie haben bis Datum X Zeit, sich zu beteiligen; danach entscheide ich, wer das Team-Wissen liefert.''
\end{itemize}

\smallskip
\textbf{5. Drei Fr\"uhwarnsignale ``Change-Programm scheitert'':}
\begin{itemize}
  \item Krankenstand im Bereich steigt $>$\,30\,\% \"uber Baseline.
  \item Anzahl der Vorschl\"age f\"ur Bot-Verbesserungen sinkt gegen null (Disengagement).
  \item Exceptions werden bewusst \"ubereskaliert (``Operator-Streik'' im Stillen).
\end{itemize}
\end{loesungbox}

\clearpage

% --- AUFGABE 8 -----------------------------------------------------------
\aufgabenkopf{8}{Fallstudie ``PharmaDistrib''}{\nivLii}{40\,min}

\ytquery{RPA pharma GxP governance audit trail exception security regulated rollout}

\begin{loesungbox}
\textbf{1. Minimum Viable Governance:}
\begin{itemize}
  \item Lifecycle (9 Phasen) verbindlich; pro Bot \emph{Phase-Sign-off} im Ticket.
  \item RACI \"uber 6 Rollen (Process Owner, Bot Owner, Developer, QA, Security, Operations).
  \item Release-Gates \emph{risikobasiert}: Low-Risk\,=\,1 Gate (QA + Standards), Medium-Risk\,=\,2 Gates (+ Process Owner), High-Risk\,=\,3 Gates (+ Security Owner). GxP-relevante Bots sind per Definition \emph{High-Risk}.
\end{itemize}

\smallskip
\textbf{2. Audit Trail Standard (GxP-tauglich):}
\begin{itemize}
  \item Felder wie in Aufgabe 5 + zus\"atzlich \texttt{GxPRelevant} (Bool), \texttt{ChangeReason}, \texttt{SignedOffBy}.
  \item Logs zentral, append-only, role-based Access (Bot Owner, QA, Audit). Operations darf \emph{lesen}, nicht \emph{l\"oschen}.
  \item Aufbewahrung 5 Jahre (Pharma-Standard), Export als signiertes Archiv f\"ur Auditor.
  \item Integrit\"atspr\"ufung w\"ochentlich (Hash-Vergleich).
\end{itemize}

\smallskip
\textbf{3. Exception Governance:}

\smallskip
\begin{tabularx}{\linewidth}{@{}p{3.0cm}p{2.4cm}Xp{1.4cm}p{2.4cm}@{}}
\toprule
\textbf{Exception} & \textbf{Owner} & \textbf{Reaktion} & \textbf{SLA} & \textbf{Eskalation} \\
\midrule
MissingBatchNumber  & Supply Chain & Queue, R\"uckfrage Werk             & 8\,h  & SC-Lead nach 24\,h \\
ERPDown             & IT-Ops       & Retry 3x; Ticket; Queue parken      & 2\,h  & Plattform nach 2\,h \\
ApprovalRequired    & Bot Owner    & Vier-Augen-Task an Genehmiger       & 24\,h & PO + Bereichsleitung \\
DuplicateInvoice    & Finance      & Stop; Pr\"ufung manuell             & 8\,h  & PO nach 24\,h \\
PortalUIChanged     & Bot Own.+Dev & Bot pausieren; Selektor-Update      & 24\,h & QA + Sec nach 48\,h \\
\bottomrule
\end{tabularx}

\smallskip
KPIs: \emph{Exception Rate} $\leq$\,12\,\%, \emph{MTTR} $<$\,4\,h; bei \"Uberschreitung Eskalation an Steuerkreis.

\smallskip
\textbf{4. Security-by-Design Guardrails:}
\begin{itemize}
  \item Bot-Accounts getrennt von Personal-Accounts; \texttt{srv\_rpa\_*}.
  \item Secrets im Vault, Rotation quartalsweise; bei Audit-Befund sofort.
  \item SoD strikt: \emph{Anlegen + Freigeben} nie im selben Account.
  \item Logging append-only, PII maskiert; Zugriff auf Logs role-based, Vier-Augen f\"ur PII-Logs.
\end{itemize}

\smallskip
\textbf{5. Change-Paket:}
\begin{itemize}
  \item \emph{Kernbotschaft:} ``Bots \"ubernehmen Routine, damit Sie Compliance und Qualit\"at sch\"arfen k\"onnen.''
  \item \emph{Rollenwandel}: Operator \textrightarrow{} Owner \textrightarrow{} Champion.
  \item \emph{Trainingsplan}: 3-stufig; verpflichtende Module zu Audit-Trail und Exception-Playbooks.
  \item \emph{Champions} je Bereich (mindestens 2 pro Bot) \textendash{} Multiplikator und Feedbackkanal.
\end{itemize}

\smallskip
\textbf{6. Priorit\"atenempfehlung (4 Wochen):}
\begin{itemize}
  \item \emph{Start:} (a) Audit-Trail-Standard + zentrales Logging, (b) Governance-Minimum (RACI + Gates), (c) drei Exception-Playbooks f\"ur Top-Bots.
  \item \emph{Verschoben:} Vollst\"andige Vendor-Audits und automatisierte Code-Quality-Checks \textendash{} sinnvoll, aber nicht in 4 Wochen umsetzbar; Auditor akzeptiert Roadmap-Commitment.
\end{itemize}
\emph{Begr\"undung:} Vor dem Audit-Sampling z\"ahlt \emph{Nachweisbarkeit} mehr als Vollst\"andigkeit \textendash{} ein sauberer Audit Trail + gelebte Governance schl\"agt einen unvollst\"andigen, aber technisch ausgefeilten Stack.
\end{loesungbox}

\clearpage

% =========================================================================
% L3 AUFGABEN
% =========================================================================
\section*{Niveau L3 \textendash{} Management \& Entscheidungsarchitektur}
\addcontentsline{toc}{section}{L3 -- Management}

% --- AUFGABE 9 -----------------------------------------------------------
\aufgabenkopf{9}{Risk Appetite \& Bot-Klassifikation}{\nivLiii}{35\,min}

\ytquery{RPA risk appetite bot classification governance high risk human in the loop}

\begin{loesungbox}
\textbf{1. Bot-Klassifikation in drei Stufen:}

\smallskip
\begin{tabularx}{\linewidth}{@{}p{3.0cm}XXX@{}}
\toprule
\textbf{Kriterium} & \textbf{Low} & \textbf{Medium} & \textbf{High} \\
\midrule
Transaktionswert         & $<$\,1.000\,\euro/Vorgang   & 1.000\,--\,50.000\,\euro       & $>$\,50.000\,\euro \\
Datensensitivit\"at      & \"offentlich / intern        & vertraulich (PII light)         & streng vertraulich (PII / Gesundheit) \\
Regulator.\ Relevanz     & keine                        & GoBD / DSGVO Standard           & GxP / KRITIS / MaRisk \\
Reversibilit\"at         & sofort r\"uckg\"angig        & in $<$\,24\,h korrigierbar      & schwer reversibel (Buchung, Zahlung, Personal) \\
\bottomrule
\end{tabularx}

\smallskip
\textbf{2. Freigabestufen je Klasse:}
\begin{itemize}
  \item \emph{Low:} 1 Gate (QA). Kein Human-in-the-Loop n\"otig.
  \item \emph{Medium:} 2 Gates (QA + Process Owner). Human-in-the-Loop bei Schwellen-\"Uberschreitung.
  \item \emph{High:} 3 Gates (QA + Process Owner + Security). Pflicht: Human-in-the-Loop f\"ur jede freigaberelevante Entscheidung; quartalsweiser Sicherheitsreview.
\end{itemize}

\smallskip
\textbf{3. Drei klassen\"ubergreifende SoD-Regeln:}
\begin{itemize}
  \item \emph{Anlegen + Freigeben} im selben Account ist verboten.
  \item \emph{\"Anderung Bankverbindung + Ausl\"osen Zahlung} sind getrennt.
  \item \emph{Login-Recht-Vergabe + Audit-Log-Zugriff} liegen nicht beim selben Bot.
\end{itemize}
Ausnahmen sind nur f\"ur \emph{Low-Risk}-Bots mit Read-only-Operationen denkbar; jede Ausnahme braucht schriftliche Vorstandsfreigabe und endet automatisch nach 6 Monaten.

\smallskip
\textbf{4. Logging \& Retention Standard:}
\begin{itemize}
  \item \emph{Low:} 1 Jahr Aufbewahrung; Standard-Logfelder.
  \item \emph{Medium:} 3 Jahre; zus\"atzlich \texttt{ChangeReason}.
  \item \emph{High:} 7 Jahre; append-only, Hash-Integrit\"atspr\"ufung w\"ochentlich, Vier-Augen f\"ur Zugriff auf PII-Logs.
\end{itemize}

\smallskip
\textbf{5. Zwei Entscheidungen unter Unsicherheit:}
\begin{itemize}
  \item \emph{Fast Track f\"ur Low-Risk-Bots erlaubt} \textendash{} Restrisiko: ein Low-Bot wird f\"alschlich als Medium klassifiziert. \emph{Fr\"uhwarnsignal:} Exception Rate $>$\,15\,\% oder Audit-Finding zu Logging-Integrit\"at \textrightarrow{} Fast Track sofort ausgesetzt, alle Low-Bots werden in 14 Tagen reklassifiziert.
  \item \emph{Human-in-the-Loop nur f\"ur High-Risk} \textendash{} Restrisiko: Medium-Bot baut Schaden auf, der h\"atte gefangen werden k\"onnen. \emph{Fr\"uhwarnsignal:} mehr als ein Incident pro Quartal in Medium-Klasse mit Schadenswert $>$\,10\,k\,\euro{} \textrightarrow{} HITL-Pflicht auf Medium erweitern.
\end{itemize}

\smallskip
\textbf{6. Vorstandsvorlage (Kernsatz):}
``Wir f\"uhren ein dreistufiges Risk-Appetite-Modell f\"ur RPA ein, das schnelle Skalierung bei \emph{Low}-Bots erm\"oglicht, \emph{Medium}-Bots strukturiert freigibt und \emph{High}-Bots ausnahmslos durch SoD, HITL und 7-Jahres-Audit-Trail absichert; die Klassifikation bindet Aufwand an Risiko und macht Verantwortlichkeit messbar \textendash{} mit klaren Stoppsignalen, falls die Wirklichkeit von der Klassifikation abweicht.''
\end{loesungbox}

\clearpage

% --- AUFGABE 10 ----------------------------------------------------------
\aufgabenkopf{10}{CoE vs.\ Federation}{\nivLiii}{30\,min}

\ytquery{RPA operating model CoE federation hub spoke automation center excellence}

\begin{loesungbox}
\textbf{1. Drei Modelle im Vergleich:}

\smallskip
\begin{tabularx}{\linewidth}{@{}p{3.0cm}XXX@{}}
\toprule
\textbf{Kriterium} & \textbf{CoE-zentral} & \textbf{F\"oderiert (mit Standards)} & \textbf{Hub-and-Spoke} \\
\midrule
Skalierungs\-geschwindigkeit & langsam (Engpass am CoE) & schnell (jeder baut) & mittel \\
Qualit\"atskonsistenz  & hoch (alle gleichen Standards) & niedrig (Standards meist papiern) & hoch, wenn Hub stark \\
Akzeptanz Fachbereich  & gering (``die da oben'')        & hoch (``mein Bot'')               & mittel-hoch \\
Engpass-Risiko         & sehr hoch (CoE-Bottleneck)      & gering, daf\"ur Wildwuchs         & hoch im Hub bei Reviews \\
Audit-Tauglichkeit     & sehr hoch                       & gering (Schatten-Bots)            & hoch, wenn Hub Sign-off macht \\
\bottomrule
\end{tabularx}

\smallskip
\textbf{2. Empfehlung} f\"ur 60 Bots, 8 Bereiche, mittlere Reife: \emph{Hub-and-Spoke (Hybrid)}. Begr\"undung: Bei dieser Gr\"o{\ss}enordnung scheitert reines CoE am Durchsatz, reine F\"oderation an Qualit\"at. Der \emph{Hub} (zentrales BPM/Security-Team) liefert Standards, Templates, Audit-Pfad und freigibt High-Risk-Bots; die \emph{Spokes} (Bereichs-Developer) bauen autonom, aber innerhalb von Leitplanken.

\smallskip
\textbf{3. Engpass-Diagnose und Hebel:}
\begin{itemize}
  \item \emph{Engpass:} Security-Review f\"ur jeden Bot.
  \item \emph{Hebel 1:} Automatisierte Code-Checks (Static Analysis) als \emph{Voraussetzung} f\"ur Security-Review \textendash{} f\"angt 60\,\% der Befunde maschinell.
  \item \emph{Hebel 2:} Templates pro Bot-Klasse (z.\,B.\ ``Read-only-Bot'' kommt ohne Vault-Review aus). Self-Service mit harten Guardrails.
  \item \emph{Hebel 3:} Reviewer-Pool aus 5 zertifizierten Bereichs-Devs (Schulung durch Hub); entlastet das Security-Team.
  \item \emph{Hebel 4:} Risk-Based Routing \textendash{} Low-Bots umgehen Security-Review komplett, Medium kriegt Sampling, High immer voll.
\end{itemize}

\smallskip
\textbf{4. Drei Fr\"uhwarnsignale ``Modell stimmt nicht mehr'':}
\begin{itemize}
  \item Anteil ``Shadow Bots'' (nicht im zentralen Repository) $>$\,15\,\%.
  \item Time-to-Production $>$\,8 Wochen f\"ur Standard-Bots.
  \item Audit-Findings zu Logging-/SoD-Bruch $>$\,3 pro Quartal.
\end{itemize}

\smallskip
\textbf{5. Faustregel:} Ein Operating Model kippt, wenn entweder \emph{Schatten-Bots} (F\"oderation au{\ss}er Kontrolle) oder \emph{Time-to-Production} (CoE-Engpass) den Hauptzweck zerst\"oren \textendash{} sp\"atestens dann braucht es einen bewussten Schwenk, nicht eine weitere Eskalations\-runde.
\end{loesungbox}

\clearpage

% --- AUFGABE 11 ----------------------------------------------------------
\aufgabenkopf{11}{Transferprojekt \textendash{} Enterprise RPA Operating Model}{\nivLiii}{50\,min}

\ytquery{enterprise RPA operating model risk appetite governance scaling automation CoE}

\begin{loesungbox}
\emph{Musterl\"osung als Entscheidungsarchitektur \textendash{} andere Konfigurationen sind m\"oglich, solange sie als Architektur (nicht als Liste) gedacht sind.}

\smallskip
\textbf{1. Risk Appetite \& Guardrails:} Drei-Klassen-Modell (Low/Medium/High) nach Wert, Sensitivit\"at, Regulatorik, Reversibilit\"at. SoD-Regeln (Anlegen+Freigeben getrennt, Bankdaten+Zahlung getrennt). Logging/Retention 1\,/\,3\,/\,7 Jahre. Human-in-the-Loop verpflichtend f\"ur High.

\smallskip
\textbf{2. Operating Model:} \emph{Hub-and-Spoke}. Hub (zentrales BPM + Security) gibt Standards, Templates, High-Risk-Sign-off; Spokes (Bereichs-Developer) bauen autonom Low/Medium. Eskalation: Spoke-Lead \textrightarrow{} Hub \textrightarrow{} CRO \textrightarrow{} Vorstand.

\smallskip
\textbf{3. Quality \& Release Management:}
\begin{itemize}
  \item \emph{Teststrategie:} Unit + Regression + Recovery; verpflichtende Negative-Tests (was, wenn ERP weg ist?).
  \item \emph{UAT:} mit echten Fachbereichsdaten in der Test-Umgebung; Sign-off Process Owner.
  \item \emph{Change-Prozess:} Standard/Normal/Emergency \textendash{} Emergency mit Post-Mortem Pflicht.
  \item \emph{Versionierung:} SemVer; Rollback-Plan vor Go-Live.
  \item \emph{Monitoring-Katalog:} Health, Queue-Tiefe, Erfolgsquote, Exception Rate, MTTR, UI-Drift-Warnung.
\end{itemize}

\smallskip
\textbf{4. Exception Governance:} Codes zentral, Playbooks pro Top-20-Exception, KPIs (Exception Rate $\leq$\,12\,\%, MTTR $<$\,4\,h, Rework $<$\,5\,\%). Quartalsweise Root-Cause-Reviews; daraus eigene Backlog-Tickets im Improve.

\smallskip
\textbf{5. Security-by-Design Blueprint:} Bot-Accounts getrennt; Vault (Rotation quartalsweise); Logging append-only mit Hash-Pr\"ufung; PII maskiert; Portal-Selektoren versioniert; Drittsysteme \"uber Service-Layer (keine direkten UI-Zugriffe wo m\"oglich).

\smallskip
\textbf{6. Change \& Enablement:}
\begin{itemize}
  \item Stakeholder-Plan: Bereichsleitung (Sponsor), Operator (Owner), QA, Security, Betriebsrat.
  \item Trainingsrollen: Operator (4\,h) \textrightarrow{} Owner (16\,h) \textrightarrow{} Champion (40\,h + Mentoring).
  \item Kommunikationsnarrativ: ``Bots nehmen Aufgaben, nicht W\"urde'' \textendash{} aber mit ehrlichem Hinweis auf m\"ogliche Stellenstrukturen.
  \item Erfolgsmessung: Akzeptanz-Survey, Skill-Shift (\% Operator $\rightarrow$ Owner), Bot-Vorschl\"age aus dem Fachbereich.
\end{itemize}

\smallskip
\textbf{7. Zwei Entscheidungen unter Unsicherheit:}
\begin{enumerate}
  \item \emph{CoE vs.\ F\"oderation:} Hub-and-Spoke gew\"ahlt. \emph{Annahme:} Standards + Templates + Sampling halten Qualit\"at; F\"oderation skaliert. \emph{Verworfene Alternative:} reines CoE \textendash{} f\"angt Backlog von $>$\,30 Bots nicht ab. \emph{Fr\"uhwarnsignal:} Shadow-Bot-Rate $>$\,15\,\% oder Incident-Rate je Bot $>$\,2/Quartal \textrightarrow{} Hub-Kontrolle versch\"arfen (z.\,B.\ Mandatory Review f\"ur alle Medium-Bots).
  \item \emph{Fast Track Grenze:} Low-Risk-Bots d\"urfen direkt nach Gate 1 in Prod. \emph{Annahme:} Read-only + niedrige Schwellen f\"angt 90\,\% des Restrisikos. \emph{Verworfene Alternative:} Fast Track f\"ur alle Klassen (zu riskant). \emph{Fr\"uhwarnsignal:} Exception Rate eines Fast-Track-Bots $>$\,15\,\% oder Audit-Finding zu Logging-Integrit\"at \textrightarrow{} Fast Track wird f\"ur 90 Tage ausgesetzt.
\end{enumerate}

\smallskip
\textbf{8. Anschluss an Strategie (drei Argumente):}
\begin{enumerate}
  \item \emph{Compliance-Risiko:} Auditf\"ahige Trails und gelebte SoD verhindern Bu{\ss}gelder und Reputations\-sch\"aden \textendash{} ein Audit-Finding kostet typischerweise mehr als ein Jahr CoE-Budget.
  \item \emph{Skalierbarkeit:} Hub-and-Spoke macht Wachstum auf 200+ Bots m\"oglich, ohne dass Security oder QA zum Bremsklotz werden.
  \item \emph{Time-to-Value:} Risk-Based Gates verk\"urzen die durchschnittliche Time-to-Production um 40\,\% gegen\"uber zentralisiertem CoE und liefern messbar fr\"uher ROI.
\end{enumerate}
\end{loesungbox}

\clearpage

\section*{Abschluss}
\thispagestyle{plain}

\noindent Die Musterl\"osungen sind eine Diskussionsbasis. Wenn Ihre eigene L\"osung anders ist, fragen Sie: \emph{Habe ich andere Annahmen \textendash{} oder eine andere Risikoabw\"agung?} Beides ist legitim, solange die Logik nachvollziehbar bleibt und Fr\"uhwarnsignale benannt sind.

\medskip
\begin{infobox}[N\"achster Schritt]
Bringen Sie Ihre L\"osungen in die Coaching-Sitzung. Drei Leitfragen: Welche Annahmen liegen Ihren Klassifikationen zugrunde? Wo haben Sie bewusst Restrisiko akzeptiert \textendash{} und welches Fr\"uhwarnsignal w\"urde Sie umsteuern lassen? Wie sichern Sie, dass Governance nicht zur B\"urokratie wird, sondern Skalierung erm\"oglicht?
\end{infobox}

\vspace{1cm}
\noindent{\color{lpAccent}\rule{\linewidth}{0.6pt}}\\[4pt]
{\footnotesize\sffamily\color{lpGray}Lernplatt \textperiodcentered{} by Can Siebert \textperiodcentered{} Kurs 71 (Dig 42) \textperiodcentered{} Tag 12 \textperiodcentered{} L\"osungsheft \textperiodcentered{} \textcopyright{} 2026}

\end{document}
